\documentclass[11pt, a4paper]{article}
\usepackage[utf8]{utf8}
\usepackage[margin=1in]{geometry}
\usepackage{hyperref}
\usepackage{enumitem}
\usepackage{titlesec}
\usepackage{xcolor}

\definecolor{nasablue}{RGB}{11, 61, 145}

\hypersetup{
    colorlinks=true,
    linkcolor=nasablue,
    urlcolor=nasablue
}

\title{\textbf{\LARGE AstraFlame: AI-Powered Microgravity Combustion Intelligence \& Spacecraft Fire Safety Dashboard}}
\author{\textbf{NASA Space Apps Challenge 2026 Submission}}
\date{\today}

\begin{document}

\maketitle

\begin{abstract}
Fire behavior in microgravity fundamentally defies terrestrial combustion intuition. In the absence of natural buoyant convection, flames transition from upward teardrop structures into spherical, diffusion-dominated regimes where airflow velocity, ambient oxygen concentration ($\mathrm{O_2\%}$), and structural material properties dictate extinction and spread. To address the NASA Space Apps Challenge (\textit{Flame in Freefall: AI-Powered Fire Safety Insights from Microgravity Combustion Data}), we present \textbf{AstraFlame}, an interactive, scientifically defensible web dashboard that summarizes, ranks, interprets, and visualizes microgravity fire dynamics to optimize fire safety protocols for human spaceflight.

AstraFlame integrates empirical data across five core NASA Physical Science Informatics (PSI) repositories: \textbf{FLEX} (droplet flammability/extinction), \textbf{BASS-II} (solid material flame spread), \textbf{SAFFIRE-I} (large-scale spacecraft burns), \textbf{SAME} (smoke aerosol generation), and \textbf{SPICE} (soot formation and flame length). The backend features a hybrid machine-learning and physics router. Machine learning models (XGBoost/Random Forest) trained on high-density datasets ($\text{FLEX}, N=274$) predict liquid fire extinction probabilities ($P(\text{extinction})$) and burn rates, while analytical physics correlations process solid fuel spread dynamics ($\text{BASS-II}, N=129$) and flame geometry ($\text{SPICE}, N=469$). Sparse flight datasets like SAFFIRE-I ($N=4$) and SAME ($N=8$) serve as ground-truth real-scale benchmarks to prevent model hallucination.

To ensure complete transparency, predictions are paired with \textbf{SHAP (SHapley Additive exPlanations)} feature attribution to quantify environmental drivers (e.g., $\mathrm{O_2\%}$ vs. airflow velocity), along with explicit domain validity checks that suppress predictions when inputs exceed NASA experiment boundaries. The frontend dashboard delivers:
\begin{enumerate}[leftmargin=*]
    \item \textbf{Interactive HTML5 Canvas Engine:} Animates spherical microgravity flames, blackbody radiation color shifts, and airflow wake stretch in real time.
    \item \textbf{Traceable Material Flammability Risk Scoreboard:} Ranks materials (e.g., PMMA, Nomex, Silicone) based on flame spread velocity and minimum extinction oxygen thresholds, linked directly to NASA PSI Investigation IDs.
    \item \textbf{What-If Ventilation Shutdown Simulator:} Demonstrates forced-airflow reduction ($0\text{ cm/s}$) to model passive flame quenching and optimize spacecraft ventilation emergency protocols.
    \item \textbf{IoT Sensor Integration:} Accepts live telemetry from an ESP32 microcontroller stream ($\mathrm{O_2}$, pressure, temperature) to evaluate real-time spacecraft cabin flammability risks.
\end{enumerate}

By uniting empirical NASA dataset provenance, calibrated machine learning, and real-time visualization, AstraFlame converts complex microgravity physics into actionable, life-saving decision support for mission control and spacecraft design engineers.
\end{abstract}

\end{document}